% Generated draft: requires mathematical and source review.
\documentclass[11pt]{article}
\usepackage[margin=25mm]{geometry}
\usepackage{fontspec}
\setmainfont{Latin Modern Roman}
\usepackage{amsmath,amssymb,mathtools,unicode-math}
\setmathfont{Latin Modern Math}
\usepackage{enumitem,needspace,xurl,hyperref}
\setlength{\parindent}{0pt}
\setlength{\parskip}{5pt}
\setlength{\emergencystretch}{3em}
\newcommand{\N}{\mathbb{N}}
\newcommand{\Z}{\mathbb{Z}}
\newcommand{\Q}{\mathbb{Q}}
\newcommand{\R}{\mathbb{R}}
\newcommand{\C}{\mathbb{C}}
\newcommand{\F}{\mathbb{F}}
\newcommand{\A}{\mathbb{A}}
\newcommand{\PP}{\mathbb{P}}
\newcommand{\E}{\mathbb{E}}
\newcommand{\eps}{\varepsilon}
\newcommand{\dd}{\,\mathrm{d}}
\begin{document}
\section*{UnsolvedMath ID: 800492}
\textbf{Source problem code:} catalogue entry
\subsection{Definitions and mathematical statement}
A parameterized family of linear systems may be written schematically as
\[
\mathcal{B}=\{A_{\theta}x=b_{\theta}:\theta\in\Theta\},
\]
where \(\Theta\) denotes the admissible model and discretization parameters, and \(A_{\theta}\) and \(b_{\theta}\) are the matrix and right-hand side produced by discretizing a partial differential equation (PDE). This notation does not prescribe a scalar field, matrix dimensions, parameter domain, PDE class, or boundary or initial conditions; these are not specified in the supplied record.

\textbf{Requested development problem.} Develop software providing a pipeline for generating benchmark parameterized linear systems from discretizations of model PDE problems. The benchmark families should be generally accepted as relevant in practice and suitable for theoretical analysis. They should permit generation of a large collection of sample systems so that the range of characteristics and solver behaviors within a problem class can be investigated.

The source motivates this focus by scientific and engineering applications and by connections between PDE discretizations and graph Laplacians studied in theoretical computer science. It contrasts a parameterized generation pipeline with collections of individual matrices, citing SuiteSparse as an example. It does not require every generated matrix to be a graph Laplacian.

\textbf{Unresolved specification.} The source gives a software and benchmark-design challenge rather than a fully quantified conjecture. It does not define criteria for practical relevance, general acceptance, suitability for theoretical analysis, or sufficient coverage of a problem class. Consequently, the supplied record alone does not determine a precise mathematical success predicate.
\subsection{Short English statement}
Develop software that generates parameterized benchmark linear systems from model PDE discretizations, with families useful in practice and amenable to theoretical analysis. The source leaves the PDEs, parameters, discretizations, and acceptance criteria unspecified.
\subsection{Sources}
\begin{itemize}
\item Local UnsolvedMath registry, permanent ID \(800492\), exact source problem code \texttt{catalogue entry}, title ``Linear Systems and Eigenvalue Problems --- Problem 1.'' The supplied record attributes the proposal to Amsel, Baumann et al.\ (2026) and labels it \texttt{open}; this label is supplied metadata, not an independently verified mathematical assessment.
\item Amsel, Baumann et al., \emph{Linear Systems and Eigenvalue Problems: Open Questions from a Simons Workshop} (2026), Problem 1. Supplied source URL: \url{https://arxiv.org/abs/2602.05394}. The record reports access on 2026-10-16 and extraction by \texttt{arxiv\_latex}; these are supplied provenance details, not a claim of independent source access or verification.
\item The supplied statement mentions the SuiteSparse matrix collection and the reference label \texttt{[DavisHu2011]}. No full bibliographic entry or URL for that reference is supplied.
\end{itemize}
\end{document}
